\section{Regular Ideals}

Let A be a k-algebra. Recall (III, §1, No.~2, p.~430) that a left ideal of A is a k-submodule \(\mathfrak{a}\) of A such that the relations \(a \in A\) , \(x \in \mathfrak{a}\) imply \(ax \in \mathfrak{a}\) . We define the notions of right ideal and two-sided ideal likewise. If A is a k-algebra and \(\mathfrak{a}\) a two-sided ideal of A, then when passing to the quotients, the multiplication in A defines a k-algebra structure on the k-module A/\(\mathfrak{a}\) (loc. cit.).

A left ideal of A is called maximal if it is a maximal element of the set of proper left ideals of A for the inclusion.

DEFINITION 1. --- Let A be a k-algebra and \(\mathfrak{a}\) a left ideal of A. An element u of A such that au - a belongs to \(\mathfrak{a}\) for every \(a \in A\) is called a right unit modulo \(\mathfrak{a}\) . We say that the ideal \(\mathfrak{a}\) is regular if there exists a right unit modulo \(\mathfrak{a}\) .

Likewise, we say that a right ideal \(\mathfrak{b}\) of A is regular if A has a left unit modulo \(\mathfrak{b}\) , that is, an element v such that \(va - a \in \mathfrak{b}\) for every \(a \in A\) . When an ideal is two-sided, one needs to specify whether it is regular as a left ideal or as a right ideal.

When the algebra A is unital, the unit element of A is a right unit modulo \(\mathfrak{a}\) for every left ideal \(\mathfrak{a}\) of A; in this case, every left (or right) ideal of A is regular. On the other hand, if A is a k-algebra whose elements are all nilpotent, then A is the only (left or right) ideal of A that is regular.

Let \(\mathfrak{a}\) be a regular left ideal of A and \(u\) a right unit modulo \(\mathfrak{a}\) . If \(\mathfrak{b}\) is a left ideal of A containing \(\mathfrak{a}\) , then \(u\) is a right unit modulo \(\mathfrak{a}\) , so \(\mathfrak{b}\) is regular. So the left ideals of A that are maximal and regular are the maximal elements of the set of regular proper left ideals of A. Moreover, a left ideal \(\mathfrak{b}\) of A containing \(\mathfrak{a}\) is proper if and only if \(u\) does not belong to \(\mathfrak{b}\) . The set of proper left ideals \(\mathfrak{b}\) of A containing \(\mathfrak{a}\) , ordered by inclusion, is therefore an inductive set. Theorem 2 of Set Theory, III, §2, No.~4, p.~154 applied to this set gives the following result.

PROPOSITION 1. --- Every regular proper left (resp. right) ideal of A is contained in a regular maximal left (resp. right) ideal.

PROPOSITION 2. --- Let A be a commutative k-algebra. An ideal \(\mathfrak{a}\) of A is a regular maximal ideal if and only if the pseudoring A/ \(\mathfrak{a}\) (I, §8, No.~1, p.~98) is a field.

An element u of A is a (right or left) unit modulo \(\mathfrak{a}\) if and only if the canonical image of u in A/\(\mathfrak{a}\) is a unit element of the algebra A/\(\mathfrak{a}\) . Therefore, the ideal \(\mathfrak{a}\) is regular if and only if the algebra A/\(\mathfrak{a}\) is unital. Suppose that these conditions are satisfied.

The mapping \(\mathfrak{b} \mapsto \mathfrak{b}/\mathfrak{a}\) is a bijection from the set of ideals of the algebra A containing \(\mathfrak{a}\) to the set of ideals of the algebra A/\(\mathfrak{a}\) . These are the ideals of the ring A/\(\mathfrak{a}\) because this algebra is unital. Finally, by Theorem 1 of I, §9, No.~1, p.~115, the ring A/\(\mathfrak{a}\) is a field if and only if it is nonzero and its only ideals are 0 and itself. The proposition follows.

Examples. --- 1) Let V be an infinite-dimensional vector space over a commutative field K. Let \(\operatorname{End}_{\mathrm{K}}^{\mathrm{f}}(\mathrm{V})\) be the K-subalgebra of \(\operatorname{End}_{\mathrm{K}}(\mathrm{V})\) consisting of the endomorphisms of finite rank. Let W be a linear subspace of V, and let \(\mathfrak{a}_{W}\) be the set of elements of \(\operatorname{End}_{\mathrm{K}}^{\mathrm{f}}(\mathrm{V})\) whose kernel contains W; it is a left ideal of \(\operatorname{End}_{\mathrm{K}}^{\mathrm{f}}(\mathrm{V})\) . An element u of \(\operatorname{End}_{\mathrm{K}}^{\mathrm{f}}(\mathrm{V})\) is a right unit modulo \(\mathfrak{a}_{W}\) if and only if we have \(u(x)=x\) for every \(x\in W\) . Such an element u exists, that is, \(\mathfrak{a}_{W}\) is regular, if and only if W is finite-dimensional. The ideal \(\mathfrak{a}_{W}\) is maximal and regular if and only if W has dimension 1.

*2) Let \(\mathrm{T}\) be a locally compact space, and let \(\mathcal{C}_0(\mathrm{T})\) be the commutative \(\mathbf{C}\) -algebra of continuous mappings from \(\mathrm{T}\) to \(\mathbf{C}\) that tend to 0 at infinity (Gen.~Top., X, §4, No.~4, p.~316); it is unital if and only if \(\mathrm{T}\) is compact. Let \(\mathrm{F}\) be a closed subset of \(\mathrm{T}\) , and let \(\mathfrak{a}_{\mathrm{F}}\) be the set of elements of \(\mathcal{C}_0(\mathrm{T})\) whose restriction to F is the zero function; it is an ideal of \(\mathcal{C}_{0}(T)\) . An element u of \(\mathcal{C}_{0}(T)\) is a unit modulo \(\mathfrak{a}_{F}\) if and only if we have \(u(t)=1\) for every \(t\in F\) . Such an element u exists, that is, \(\mathfrak{a}_{F}\) is regular, if and only if F is compact. The mapping \(t\mapsto \mathfrak{a}_{\{t\}}\) is a bijection from T to the set of regular maximal ideals of \(\mathcal{C}_{0}(T)\) (TS, I, §3, n° 2, p.~32, corollaire 1). Suppose that T is not compact, and denote by \(\mathfrak{a}\) the subset of \(\mathcal{C}_{0}(T)\) consisting of the functions with compact support; then \(\mathfrak{a}\) is an ideal of \(\mathcal{C}_{0}(T)\) that is not contained in any regular ideal of \(\mathcal{C}_{0}(T)\) .

\begin{enumerate}
\def\labelenumi{\arabic{enumi})}
\setcounter{enumi}{2}
\tightlist
\item
  Let \(L^{1}(\mathbf{R})\) be the convolution algebra of the locally compact group R. Recall (cf.~Int., VIII, §4, No.~5, p.~38) that \(L^{1}(\mathbf{R})\) is the space of classes of functions on R that are integrable for the Lebesgue measure. The product of the classes of two functions f and g is the class of the function \(f * g\) defined by the formula
\end{enumerate}

\[
(f * g) (s) = \int_ {- \infty} ^ {+ \infty} f (t) g (s - t) d t
\]

for almost every \(s \in \mathbf{R}\) . The algebra \(\mathrm{L}^1(\mathbf{R})\) is not unital. For any \(a\) in \(\mathbf{R}\) , denote by \(\mathfrak{m}_a\) the set of elements \(f\) of \(\mathrm{L}^1(\mathbf{R})\) satisfying

\[
\int_ {- \infty} ^ {+ \infty} f (t) e ^ {- i a t} d t = 0.
\]

By TS, II, §3, n° 2, p.~252, théorème 1, the mapping \(a \mapsto \mathfrak{m}_a\) is a bijection from \(\mathbf{R}\) to the set of regular maximal ideals of the algebra \(\mathrm{L}^1(\mathbf{R})\) .*

